% Capitulo: Dados para treinar IA
\chapter{Dados para treinar IA}
\label{cap:treinar-ia}

Quando o consumidor do dado é um modelo que vai aprender com ele, a modelagem muda de
natureza. Deixa de ser sobre tabelas e métricas e passa a ser sobre exemplos: quais entram,
como são rotulados, como são versionados e documentados, e o que nunca pode entrar. Este
capítulo trata dessas decisões. Ele parte de uma constatação prática: a maioria das empresas
não vai treinar um modelo de linguagem do zero, mas quase todas vão ajustar modelos
existentes ou, no mínimo, precisar avaliar se a IA que usam responde certo. As três
coisas pedem dados organizados com cuidado.

\section{Quando o consumidor é um modelo}
\label{sec:consumidor-modelo}

Há três necessidades de dados diferentes quando se fala em treinar IA, e elas aparecem nas
empresas em ordem crescente de frequência.

A primeira é treinar um modelo do zero. Para modelos de linguagem de uso geral, isso exige
volumes de texto e de processamento que estão fora do alcance de quase todas as empresas.
Para modelos menores e específicos, como um classificador de documentos ou um detector de
anomalias, continua sendo comum, e os cuidados do capítulo~\ref{cap:modelagem-cd} valem
integralmente.

A segunda é ajustar um modelo existente (\emph{fine-tuning}): partir de um modelo pronto e
treiná-lo um pouco mais com exemplos da empresa, para que aprenda um vocabulário, um formato
de resposta ou uma tarefa específica. Aqui, poucas centenas ou milhares de exemplos bons
podem fazer muita diferença, e a qualidade de cada exemplo pesa mais do que a quantidade.

A terceira, e mais comum, é avaliar. Toda empresa que usa IA para responder perguntas,
classificar documentos ou gerar consultas precisa saber se ela acerta, e para isso precisa de
um conjunto de avaliação: perguntas ou tarefas com a resposta certa conhecida. Esse conjunto
é, na prática, um dado de treino que todo time de dados vai precisar montar, mesmo sem
treinar nada. Ele é a forma de saber se uma mudança no modelo, no contexto ou nos dados
melhorou ou piorou as respostas, e o capítulo~\ref{cap:contexto-ia} volta a ele.

\section{Curadoria, rotulagem e versionamento}
\label{sec:curadoria}

Em dados para treino, poucos exemplos bons valem mais do que muitos ruins. Exemplos
duplicados fazem o modelo decorar em vez de generalizar. Exemplos contraditórios, em que a
mesma situação tem respostas diferentes, ensinam o modelo a ser inconsistente. Exemplos fora
do domínio desviam o aprendizado. A curadoria (selecionar, limpar e equilibrar os exemplos)
é, por isso, a parte mais importante do trabalho, e a mais parecida com o trabalho de
qualidade de dados que os capítulos anteriores descreveram.

Muitos exemplos precisam de rótulo: a categoria correta de um documento, a resposta certa de
uma pergunta, se uma consulta gerada está correta ou não. Rotular é um trabalho humano, e
envolve decisões que precisam ser explícitas. Quem rotula? Com que critério escrito? O que
fazer nos casos ambíguos? E como saber se os rotuladores concordam entre si? Para essa
última pergunta há medidas clássicas de concordância, que descontam o acordo que
aconteceria por acaso \parencite{cohen1960}. Se dois especialistas do negócio discordam em
um terço dos casos, o problema não é o modelo: é que a regra não está clara, e nenhum
modelo vai aprender uma regra que as próprias pessoas não aplicam de forma consistente.

O conjunto de dados precisa ser versionado como se versiona código (capítulo~\ref{cap:dados-codigo}).
Cada modelo treinado ou avaliado deve estar associado a uma versão exata do conjunto, para
que seja possível reproduzir o resultado, comparar versões e entender por que um modelo novo
se comporta diferente do anterior.

E a separação entre treino, validação e teste deve acontecer antes de qualquer outra coisa,
e ser respeitada. Se exemplos do teste aparecem no treino, ou se o teste é usado
repetidamente para escolher entre versões, a avaliação deixa de medir o que o modelo faria
com dados novos. É uma forma de vazamento, parente da que o capítulo anterior descreveu
(\veja{sec:vazamento}).

\section{Documentar o conjunto de dados}
\label{sec:datasheets}

Um conjunto de dados para treino sem documentação é um risco: ninguém sabe de onde veio, o
que contém, para que serve e para que não deve ser usado. Em 2018, um grupo de
pesquisadores propôs que todo conjunto de dados viesse acompanhado de uma ficha técnica, nos
moldes das fichas de componentes eletrônicos, respondendo a perguntas sobre motivação,
composição, processo de coleta, pré-processamento, usos recomendados, distribuição e
manutenção \parencite{gebru2021}. Uma proposta parecida, para os próprios modelos, descreve
o desempenho de cada modelo em grupos diferentes e as situações em que ele não deve ser
usado \parencite{mitchell2019}.

Para um time de dados, a ficha não precisa ser longa. Algumas perguntas já evitam a maior
parte dos problemas: por que este conjunto foi criado, de quais fontes e em que período veio,
quem rotulou e com que critério, que limpezas foram feitas, o que ele não cobre, e quem é o
responsável por ele. É o equivalente, para dados de treino, do que a documentação do modelo
dimensional sempre foi para o BI: o que permite que outra pessoa use o dado sem precisar
perguntar a quem o fez.

\section{Privacidade, direitos e viés}
\label{sec:privacidade-vies}

Dado pessoal que entra num treino é difícil de tirar depois, porque o modelo pode ter
aprendido algo a partir dele. Por isso a limpeza precisa acontecer antes. A Lei Geral de
Proteção de Dados distingue o dado anonimizado, que deixa de ser considerado dado pessoal
quando a identificação não é mais possível por meios razoáveis, do dado pseudonimizado, em
que a identificação continua possível com uma informação adicional guardada à parte
\parencite{lgpd2018}. Na prática, isso significa que trocar nomes por códigos não é
suficiente para tratar o dado como não pessoal. \pendente{conferir os artigos da LGPD sobre
anonimização e pseudonimização e citar os números}

Além da privacidade, há os direitos sobre o conteúdo. Documentos internos, textos de
clientes, conteúdo de terceiros e dados comprados têm termos de uso, e nem todos permitem o
uso para treinar ou ajustar modelos. Vale verificar antes, com a área jurídica, e registrar a
verificação na ficha do conjunto.

Por fim, a representatividade. Um conjunto que só tem exemplos de uma região, de um tipo de
produto ou de um perfil de cliente ensina o modelo a acertar nesses casos e a errar nos
outros. Um classificador de solicitações treinado só com dados da matriz pode ir mal nas
filiais, onde o vocabulário e os problemas são outros. A ficha técnica deve dizer o que o
conjunto cobre, e a avaliação deve medir o desempenho separadamente em cada grupo relevante,
e não só na média.

\begin{destaque}
Os modelos de IA são, em grande parte, aquilo com que foram treinados e avaliados. Cuidar
desses dados é cuidar do comportamento do modelo.
\end{destaque}

\rascunho{
  Relato opcional: uma experiência de montar um conjunto de avaliação ou de ajuste para
  uma IA usada no trabalho, e o que deu trabalho (rotulagem, casos ambíguos, privacidade).
}

\begin{leituras}
  \leitura{gebru2021}{A proposta de documentar cada conjunto de dados como se documenta
    um componente eletrônico: origem, composição, usos recomendados e limites.}
  \leitura{mitchell2019}{A proposta equivalente para modelos, com atenção ao desempenho em
    grupos diferentes.}
  \leitura{huyen2022}{Os capítulos sobre dados de treino tratam amostragem, rotulagem e
    desbalanceamento de forma prática.}
\end{leituras}
